\section{The unique neutral limiting profile}
\label{prof:section}

Fix one datum in each retained set $G_\ell$ of
Proposition~\ref{scales:good}.  The resulting fields are deterministic;
we use its notation and all its bounds below.  In particular,
$a_\ell\to0$, $q_\ell\to\infty$, and the exact neutral mass identity
\eqref{scales:mass-field} holds, even though
$\lambda_\ell=Z_\ell s_\ell^3\to\infty$ by admissibility.

We first obtain compactness of the signed fields from neutrality
(Lemma~\ref{prof:compactness}), then transfer the positive barrier
(Lemma~\ref{prof:lower}), and finally identify the limit by dilation
(Lemma~\ref{prof:rigidity}).
We do not assume that $F_\ell$ is pointwise nonnegative.

\begin{lemma}[Compactness from the neutral spherical mean]
\label{prof:compactness}
The sequence $(F_\ell)$ is locally uniformly precompact on
$\R^3\setminus\{0\}$.  Every locally uniform subsequential limit $F$
is locally Lipschitz and satisfies
\begin{equation}
\label{prof:limit-equation}
 \Delta F=f(F),\qquad F(x)\le C_{\rm cap}|x|^{-4},\qquad
 \frac1{4\pi}\int_{\mathbb S^2}F(d\omega)\,\dd\omega\ge0
 \quad(d>0).
\end{equation}
Along the same subsequence,
$\sigma_\ell\to k(F_+)^{3/2}$ locally uniformly away from zero.
Moreover, there is a constant $C_*$, depending only on $C_{\rm cap}$
and $k$, such that
\begin{equation}
\label{prof:absolute}
 |F(x)|\le C_*|x|^{-4}\qquad(x\ne0).
\end{equation}
The family of dilates $\{D^4F(D\,\cdot):D>0\}$ is also locally
uniformly precompact away from zero.
\end{lemma}

\begin{proof}
Newton's spherical-average identity is
\[
 \frac1{4\pi}\int_{\mathbb S^2}\frac{\dd\omega}{|d\omega-z|}
 =\frac1{\max(d,|z|)}.
\]
It follows either by direct integration in the angle with $z$, or by
the mean-value property inside and outside the sphere.  Applying it
to \eqref{scales:mass-field} and using the exact equality of the nuclear
coefficient and total density mass gives
\begin{equation}
\label{prof:mean}
 \overline F_\ell(d):=
 \frac1{4\pi}\int_{\mathbb S^2}F_\ell(d\omega)\,\dd\omega
 =\int_{|z|>d}\left(\frac1d-\frac1{|z|}\right)
                    \sigma_\ell(z)\,\dd z\ge0.
\end{equation}
All these integrals are finite for each $\ell$.  No upper bound on
$\lambda_\ell$ is needed.

Fix an annulus $\alpha\le |x|\le\beta$ with $0<\alpha<\beta<\infty$.
For large $\ell$, it lies in the band of \eqref{scales:band}.  The cap
and \eqref{prof:mean} imply, for every radius in this annulus,
\begin{equation}
\label{prof:sphere-lone}
 \int_{\mathbb S^2}|F_\ell(d\omega)|\,\dd\omega
 =2\int_{\mathbb S^2}(F_\ell(d\omega))_+\,\dd\omega
    -\int_{\mathbb S^2}F_\ell(d\omega)\,\dd\omega
 \le 8\pi C_{\rm cap}d^{-4}.
\end{equation}
Radial integration therefore bounds the full $L^1$ norm of $F_\ell$
on the annulus.  Equation~\eqref{scales:band} also bounds
$\sigma_\ell$ uniformly there.

Here is the local estimate we shall use twice.  Take a ball $Q$ whose
closure stays away from zero, and a larger fixed annulus containing
its closure.  The localized potential
\[
 P_\ell=K*(\sigma_\ell\1_Q)
\]
has uniformly bounded values and first derivatives: this follows
from the density bound and the integrability of $|z|^{-1}$ and
$|z|^{-2}$ on bounded subsets of $\R^3$.  Since
$\Delta F_\ell=4\pi\sigma_\ell$ on $Q$, the function
$h_\ell=F_\ell+P_\ell$ is distributionally harmonic there.  Its
$L^1(Q)$ norm is bounded by \eqref{prof:sphere-lone} and the bound on
$P_\ell$.  On every smaller concentric ball, harmonic interior
estimates bound $h_\ell$ and its gradient by this $L^1$ norm.  For
completeness, convolution with a fixed smooth radial averaging kernel
reproduces a harmonic function in the interior; differentiating that
convolution gives these bounds.  The same argument applies to
continuous distributionally harmonic functions, by first mollifying
on interior sets.  Subtracting $P_\ell$ gives uniform local value and
Lipschitz bounds for $F_\ell$.

A covering by such balls, followed by a diagonal
Arzel\`a--Ascoli argument, gives local uniform precompactness and a
locally Lipschitz limit.  The error in \eqref{scales:band} tends
uniformly to zero on every fixed compact annulus, so
$\sigma_\ell\to k(F_+)^{3/2}$ there.  Passing to distributions, to the
cap, and to the spherical averages proves
\eqref{prof:limit-equation}.

To compare with $U_\ell$ on large spheres, we next strengthen the
one-sided cap to an absolute bound for $F$.
For $D>0$ set $F^{(D)}(x)=D^4F(Dx)$.  The equation, cap, and
nonnegative spherical means in \eqref{prof:limit-equation} are
invariant under this dilation.  In particular, on a fixed annulus
the same argument as in \eqref{prof:sphere-lone} bounds the $L^1$
norm of every $F^{(D)}$, while the source
$k(F^{(D)}_+)^{3/2}$ is uniformly bounded by the cap.  The local
estimate just proved is therefore uniform in $D$.  It yields both
precompactness of the dilates and a uniform bound
$|F^{(D)}(\omega)|\le C_*$ for $|\omega|=1$.  Writing $x=D\omega$
proves \eqref{prof:absolute}.
\end{proof}

\begin{lemma}[Transfer of the positive barrier]
\label{prof:lower}
Every limit $F$ in Lemma~\ref{prof:compactness} satisfies
\begin{equation}
\label{prof:positive}
 B|x|^{-4}\le F(x)\le C_{\rm cap}|x|^{-4}\qquad(x\ne0).
\end{equation}
\end{lemma}

\begin{proof}
We compare the barrier and the screened field on a fixed ball after
subtracting a Coulomb correction that vanishes away from the origin.
Work along the subsequence converging to $F$.  Fix $S>0$; eventually
$a_\ell<S<q_\ell$.  Define
\begin{equation}
\label{prof:correction}
 e_\ell(x)=q_\ell^{-8}|x|^{-6}
                    \1_{\{a_\ell\le|x|\le S\}},
 \qquad L_\ell=K*(p_\ell+e_\ell)\ge0.
\end{equation}
Subtracting the Poisson equation for $F_\ell$ from
\eqref{scales:invariant} cancels the nucleus and the entire inner
density source.  Since
$4\pi\sigma_\ell\le f(F_\ell)+4\pi e_\ell$ on
$a_\ell\le|x|<S$ and $\Delta(-L_\ell)=4\pi(p_\ell+e_\ell)$, the
result is
\begin{equation}
\label{prof:comparison-inequality}
 \Delta(U_\ell-F_\ell-L_\ell)
 \ge \1_{\{|x|\ge a_\ell\}}
                  \bigl(f(U_\ell)-f(F_\ell)\bigr)
 \quad\hbox{in }B(0,S).
\end{equation}
This is an inequality on the whole ball, including the origin and
the sphere $|x|=a_\ell$; no restriction or differentiation across an
interface has been made.

Put
\[
 c_S=(C_{\rm inv}+C_*+1)S^{-4}.
\]
Uniform convergence on $|x|=S$, the bound
\eqref{prof:absolute}, the upper barrier in
\eqref{scales:barriers}, and $L_\ell\ge0$ give, for large $\ell$,
\[
 U_\ell-F_\ell-L_\ell<c_S\qquad\hbox{on }|x|=S.
\]
We justify the comparison throughout the ball.  The difference
\[
 W_\ell=U_\ell-F_\ell-L_\ell-c_S
       =(U_\ell-\lambda_\ell K)+K*\sigma_\ell-L_\ell-c_S
\]
extends continuously and locally in $H^1$ across zero.  Indeed, the
first term is locally Lipschitz by hypothesis.  The other potentials
have bounded integrable densities for each fixed $\ell$, and thus
locally bounded gradients.  In the case of $L_\ell$ the density is
also compactly supported.  These bounds need not be uniform in
$\ell$.

The strict boundary inequality implies that $(W_\ell)_+$ has compact
support in $B(0,S)$ and belongs to $H^1_0(B(0,S))$.  It may therefore
be used in \eqref{prof:comparison-inequality}, by approximation with
nonnegative smooth tests in $H^1$.  The right side is bounded on
this fixed ball for this fixed $\ell$, since its reaction is cut off
at the positive radius $a_\ell$.  On the positive set of $W_\ell$
outside that radius, we have $U_\ell>F_\ell$, because
$L_\ell+c_S>0$.  Monotonicity of $f$ consequently gives
\[
 -\int_{B(0,S)}|\nabla(W_\ell)_+|^2\,\dd x
 \ge\int_{B(0,S)}\1_{\{|x|\ge a_\ell\}}
       \bigl(f(U_\ell)-f(F_\ell)\bigr)(W_\ell)_+\,\dd x
 \ge0.
\]
It follows that $(W_\ell)_+=0$, or equivalently
\begin{equation}
\label{prof:comparison}
 U_\ell\le F_\ell+L_\ell+c_S\qquad\hbox{in }B(0,S).
\end{equation}

For each fixed $x$ with $0<|x|<S$, we next show that
$L_\ell(x)\to0$.  For large $\ell$, the support and mass bound in
\eqref{scales:error} imply
\[
 (K*p_\ell)(x)\le \frac{2s_\ell^3}{|x|}\longrightarrow0.
\]
The other source has small total mass:
\[
 \int e_\ell\,\dd x
 =\frac{4\pi}{3}q_\ell^{-8}(a_\ell^{-3}-S^{-3})
 \le\frac{32\pi}{3}q_\ell^{-5}\longrightarrow0.
\]
On the fixed ball $B(x,|x|/4)$ its density is at most
$C_xq_\ell^{-8}$.  Integrability of the Coulomb kernel bounds its
potential contribution there by $C_xq_\ell^{-8}$; on the complement
the kernel is at most $4/|x|$, so the contribution is at most
$(4/|x|)\int e_\ell$.  Hence $L_\ell(x)\to0$.

Now \eqref{prof:comparison} and the lower barrier give
\[
 F(x)+c_S\ge B|x|^{-4}.
\]
The argument holds for every fixed $S>|x|$ along the same convergent
subsequence.  Letting $S\to\infty$ proves the lower bound in
\eqref{prof:positive}; the upper bound was already known.
\end{proof}

The next rigidity statement also follows from Solovej's nonradial
Sommerfeld estimate \cite[Lemma~4.4 and Remark~4.5]{Solovej2003}:
in that estimate, take inner radius zero and use the positive
two-sided $|x|^{-4}$ bounds.  We retain a direct dilation proof to
make the precise compactness and contact argument explicit.

\begin{lemma}[Rigidity of a positive comparable profile]
\label{prof:rigidity}
Let $F$ be a continuous distributional solution of
$\Delta F=4\pi kF^{3/2}$ on $\R^3\setminus\{0\}$, and suppose
\[
 B|x|^{-4}\le F(x)\le C|x|^{-4}\qquad(x\ne0)
\]
for constants $0<B\le C<\infty$.  Then
\begin{equation}
\label{prof:homogeneous}
 F(x)=A_*|x|^{-4},\qquad
 A_*=\left(\frac{12}{4\pi k}\right)^2=\frac{81\pi^2}{8}.
\end{equation}
No radiality assumption is required.
\end{lemma}

\begin{proof}
The local potential and harmonic estimates in
Lemma~\ref{prof:compactness} apply uniformly to the dilates
$F^{(D)}(x)=D^4F(Dx)$: their values and their sources are bounded on
every fixed compact annulus.  Thus the dilates are locally uniformly
precompact, with locally Lipschitz limits satisfying the same
equation.

Write
\[
 b=\sup_{x\ne0}|x|^4F(x),\qquad
 h=\inf_{x\ne0}|x|^4F(x),
 \qquad 0<h\le b<\infty.
\]
Choose $x_n$ with $|x_n|^4F(x_n)\to b$, and let
$D_n=|x_n|$, $\omega_n=x_n/|x_n|$.  After passing to a subsequence,
$F^{(D_n)}\to v$ locally uniformly and $\omega_n\to\omega\in\mathbb S^2$.
Local equicontinuity gives
\[
 v(x)\le b|x|^{-4}\quad(x\ne0),\qquad v(\omega)=b.
\]
The equation passes to the limit because its reaction is continuous
and locally uniformly bounded.  Since
$\Delta |x|^{-4}=12|x|^{-6}$ in three dimensions, this contact implies
\begin{equation}
\label{prof:upper-contact}
 4\pi k b^{3/2}\le12b.
\end{equation}
Here is a weak proof of the contact assertion.  If the inequality
failed, the continuous function $g=v-b|x|^{-4}$ would satisfy
$g\le0$, $g(\omega)=0$, and $\Delta g\ge\delta>0$ on a sufficiently
small ball about $\omega$.  For $0<\eta<\delta/6$, the function
$g-\eta|x-\omega|^2$ is weakly subharmonic there, equals zero at the
center, and is strictly negative on the boundary.  Choose a level
strictly between its boundary maximum and zero.  Testing its
Laplacian with the positive part above that level, as in the proof of
Lemma~\ref{prof:lower}, contradicts the value at the center.
This proves \eqref{prof:upper-contact} without a classical second
derivative assumption.

Dilating instead along points approaching the infimum gives a limit
$v\ge h|x|^{-4}$ with equality at a unit point.  If
$4\pi k h^{3/2}<12h$, the same argument applies to
$h|x|^{-4}-v$, whose Laplacian is then strictly positive near its zero
maximum.  Consequently
\[
 4\pi k h^{3/2}\ge12h.
\]
Together with \eqref{prof:upper-contact}, this gives
$b\le A_*\le h$.  Since $h\le b$, both extrema equal $A_*$, proving
\eqref{prof:homogeneous}.  The argument allows $D_n$ to tend to zero,
to infinity, or to neither; no extremum of the original field had to
be attained.
\end{proof}

\begin{proposition}[Convergence of the selected fields]
\label{prof:limit}
For every choice of one retained datum at each index,
\[
 F_\ell(x)\longrightarrow A_*|x|^{-4},\qquad
 \sigma_\ell(x)\longrightarrow
 \rho_*(x):=kA_*^{3/2}|x|^{-6}
\]
locally uniformly on $\R^3\setminus\{0\}$.
\end{proposition}

\begin{proof}
Lemma~\ref{prof:compactness} gives a locally uniformly convergent
subsequence of every subsequence.  Its limit satisfies the positive
bounds of Lemma~\ref{prof:lower}, so Lemma~\ref{prof:rigidity}
identifies it as $A_*|x|^{-4}$.  Uniqueness of every subsequential
limit proves convergence of the full sequence.  The density
convergence follows from \eqref{scales:band}.
\end{proof}

\begin{corollary}[Uniformity on the retained data]
\label{prof:uniform}
For each fixed compact annulus $\mathcal A\subset\R^3\setminus\{0\}$,
\[
 \sup_{\theta\in G_\ell}
 \left\|\sigma_\ell(\theta,\,\cdot)-\rho_*\right\|_{C(\mathcal A)}
 \longrightarrow0.
\]
The corresponding assertion holds for
$F_\ell(\theta,\,\cdot)-A_*|\cdot|^{-4}$.
\end{corollary}

\begin{proof}
If either assertion failed, there would be a positive discrepancy
along a subsequence and a choice of one datum from each corresponding
$G_\ell$ realizing at least half that discrepancy.  Those selected
fields satisfy all the same deterministic hypotheses, contradicting
Proposition~\ref{prof:limit}.  This argument uses no measurable
selection: uniformity is a deterministic assertion about all the
retained data, after the conditional versions have been fixed.
\end{proof}
